\section{Sensitivity recurrences}
\label{sec:recurrences}

The joint profiles bound simultaneous changes of a fixed pair of predicates
on the same child. The following estimates hold for every input, including
inputs far from the blocks used in Lemma~\ref{lem:blocks}.

\begin{proposition}\label{prop:recurrences}
For $1\leq \ell\leq d$, write
$S'_b=S_{\ell-1,b}$ and $J'_b=J_{\ell-1,b}$.  Then
\begin{equation}\label{eq:recurrences}
\begin{aligned}
 S_{\ell,1}&\leq M^2S'_0+rS'_1,
 &\qquad J_{\ell,1}&\leq M^2J'_0+rS'_1,\\
 S_{\ell,0}&\leq tS'_0+S'_1+3J'_1,
 &\qquad J_{\ell,0}&\leq tS'_0+3J'_1.
\end{aligned}
\end{equation}
\end{proposition}

\begin{proof}
Fix $\ell$, put $H=H_\ell$ and $Q=H+1$, and recall
Equation~\eqref{eq:predicate-definition}.  For predicate index $q$, every
target tests $P_{\ell-1,Q}=1$ and every gate tests
$P_{\ell-1,Q-q}=0$ on its specified child.  The targets of clause $i$
lie in row $i$, and its gates lie in distinct other rows.  Thus each
condition of a fixed clause uses a different child.

All flips below are flips of individual coordinates of the original
Boolean input.  Such a flip affects exactly one child.  It may change
several nested predicates of that child; we make no restriction on the
number of thresholds it crosses.  Nevertheless, at a fixed parent index
it can change at most one condition of any fixed clause.  A child may
occur in several different clauses, so when counting possible flips
over clauses we use a union bound.

\emph{Transitions from one.}
Fix an input where $P_{\ell,q}=1$, and choose a clause satisfied there.
Every flip taking this predicate to zero must break that fixed clause.
It either changes a gate predicate from zero to one or changes a target
predicate from one to zero.  There are $M^2$ gate children and $r$
target children, giving at most $M^2S'_0+rS'_1$ such flips.

Next fix $q<q'$ and an input where both parent predicates are one.
Choose a clause initially satisfied for $q'$; nesting makes it satisfied
for $q$ as well.  This same clause is fixed for all flips being counted.
A flip taking both parent predicates to zero must either break one of
its targets, accounting for at most $rS'_1$ flips, or leave all its
targets true and fail a gate even for $q$.  Put
\[
 g=Q-q,\qquad g'=Q-q',\qquad 1\leq g'<g<Q.
\]
At such a gate, initially $P_{\ell-1,g'}=0$ because the clause holds
for $q'$, and nesting gives $P_{\ell-1,g}=0$.  Failure of the gate for
$q$ after the flip means $P_{\ell-1,g}=1$, and nesting then gives
$P_{\ell-1,g'}=1$.  The two distinct child predicates therefore both
change from zero to one.  Each of the $M^2$ gates contributes at most
$J'_0$ such flips.  This proves both inequalities on the one side.

\emph{The initial candidate clauses at a zero.}
Fix $q$ and an input $x$ with $P_{\ell,q}(x)=0$.  Let $\mathcal T$
be the set of clauses that at $x$ fail exactly one target and no gate,
and let $\mathcal G$ be the set that fail exactly one gate and no
target.  These sets are determined by $x$ and $q$, before a flipped
coordinate is chosen.  Every flip making $P_{\ell,q}$ one satisfies
some previously false clause.  Since only one condition of that clause
can change, the clause belongs to $\mathcal T\cup\mathcal G$, and
its unique failed condition must be repaired.

For each $j\in\mathcal T$, let $m_j\in[r]$ be the position of its
unique failed target.  If $i,j\in\mathcal T$ and $i\to j$, clause
$i$ has a passing gate in child $(j,a(i,j))$.  Its gate predicate
$P_{\ell-1,Q-q}$ is zero, so by nesting its target predicate
$P_{\ell-1,Q}$ is also zero.  The unique failed target in row $j$
is at position $m_j$, forcing $a(i,j)=m_j$.  Consequently a set of
$t$ vertices in $\mathcal T$ would violate Lemma~\ref{lem:labelling}.
Thus $|\mathcal T|<t$.  Repairing a clause in this list requires
changing its unique failed target from zero to one, and accounts for
at most $S'_0$ flips per clause.  Their total contribution is at most
$tS'_0$.

All targets in the row of a vertex of $\mathcal G$ are initially
true.  Hence an edge $i\to j$ within $\mathcal G$ gives a failed
gate of clause $i$: the gate child in row $j$ has
$P_{\ell-1,Q}=1$, which forces $P_{\ell-1,Q-q}=1$.  A clause in
$\mathcal G$ has just one failed gate, so every outdegree in the
induced tournament on $\mathcal G$ is at most one.  With
$m=|\mathcal G|$, counting its edges gives
\[
 \frac{m(m-1)}2\leq m,
 \qquad\text{and hence}\qquad m\leq3.
\]
There is at most one vertex of outdegree zero in this induced
tournament, since the edge between any two vertices leaves one of
them.

If $i\in\mathcal G$ has an outgoing edge $i\to j$ inside
$\mathcal G$, that edge is its unique failed gate.  On the
corresponding child, initially
\[
 P_{\ell-1,Q-q}=P_{\ell-1,Q}=1.
\]
Repairing the gate sets $P_{\ell-1,Q-q}=0$ and therefore also sets
$P_{\ell-1,Q}=0$.  The indices $Q-q$ and $Q$ are distinct because
$q\geq1$, so at most $J'_1$ flips repair this clause.  The fact that
such a flip destroys a target in row $j$ causes no problem: the
accepting witness is clause $i$, whose own targets and other gates
are unchanged.  The possible vertex with no outgoing edge inside
$\mathcal G$ contributes at most $S'_1$, by counting flips of its
unique failed gate.  Since $|\mathcal G|\leq3$, its entire
contribution is at most $S'_1+3J'_1$.  Together with the target
contribution this proves
$S_{\ell,0}\leq tS'_0+S'_1+3J'_1$.

\emph{Joint transitions from zero.}
Fix $q<q'$ and an input $x$ where both parent predicates are zero,
and form the same initial lists $\mathcal T,\mathcal G$ for the
weaker index $q$.  Consider a flip making both predicates one, and
choose any clause satisfied afterward for $q'$.  It is also satisfied
afterward for $q$, so it belongs to those fixed initial lists.  The
choice of this final witness may depend on the flipped coordinate;
the lists themselves do not.  A genuine final witness also had
exactly one failed condition at $q'$ initially, by the same
distinct-child argument.  Members of the weaker lists that have
additional unrepaired failures at $q'$ simply contribute no flips.

A witness in $\mathcal T$ is counted by the preceding target bound
$tS'_0$.  For a witness in $\mathcal G$, let $g=Q-q$ and
$g'=Q-q'$, so $1\leq g'<g<Q$.  Its unique failed gate at index $q$
has $P_{\ell-1,g}=1$ initially, and nesting gives
$P_{\ell-1,g'}=1$ initially as well.  To satisfy the clause for
$q'$ afterward, the same child must have $P_{\ell-1,g'}=0$, which
forces $P_{\ell-1,g}=0$.  Thus both distinct child predicates change
from one to zero, and each such clause contributes at most $J'_1$.
This reasoning applies also to the possible vertex with no outgoing
edge inside $\mathcal G$: the second parent index supplies the
second child threshold, without requiring any initial target value
at the failed gate.  The gate contribution is therefore at most
$3J'_1$.

Every joint-success coordinate is thus counted by a clause in the fixed
lists $\mathcal T$ or $\mathcal G$, so a union bound gives
$J_{\ell,0}\leq tS'_0+3J'_1$.

All bounds hold for every relevant input and for each fixed index or
pair of indices, so taking the maxima gives the proposition.  If
$H=1$, there is no pair of distinct parent indices and both joint
inequalities hold by the empty-maximum convention.
\end{proof}
